\section{The machine: a computer read off the geometry}
\label{sec:machine}

The same forty points that carry $E_8$ and a light cone also describe, with no design
choices, a quantum computer: an instruction set, a memory, a network and a supply of the
non-classical resource that makes quantum computation powerful. The Holonet papers develop
this at length; here we give the parts that are exact, and mark carefully what has and has
not been built.

\begin{figure}[htbp]
\centering
\begin{tikzpicture}[>={Stealth[length=2.2mm]},
   lay/.style={draw=#1,fill=#1!10,rounded corners=4pt,align=center,font=\sffamily\scriptsize,minimum width=27mm,minimum height=12mm}]
  \def\R{3.1}
  \foreach \nm/\c/\txt [count=\i from 0] in
     {car/ink/{\textbf{carrier}\\one photon\\time $\otimes$ frequency},
      isa/teal/{\textbf{instructions}\\$81$ Pauli frames\\$4$ opcodes, $2$ bits},
      mem/sky/{\textbf{memory}\\Steinberg $81$\\$[[240,81,3]]_3$},
      net/gold/{\textbf{network}\\$540$ cube charts\\$\le 8$ hops/digit},
      fuel/rose/{\textbf{fuel}\\$36$ magic rays\\of $40$},
      cub/sage/{\textbf{universal gate}\\the $E_6$ cubic\\as a native tick}}{
     \pgfmathsetmacro{\a}{90-60*\i}
     \node[lay=\c] (\nm) at (\a:\R) {\txt};}
  \foreach \x/\y in {car/isa,isa/mem,mem/net,net/fuel,fuel/cub,cub/car}{
     \draw[->,ink!45,line width=0.9pt] (\x) to[bend left=12] (\y);}
  \node[circle,fill=ink,text=white,font=\sffamily\small\bfseries,minimum size=19mm,align=center] at (0,0) {$\W$\\[-1pt]\scriptsize $40$ points};
\end{tikzpicture}
\caption{The Holonet stack. Every layer is a structure of $\W$ read as engineering: the
Pauli frames are the $81$ points of $\F_3^4$, the memory is the Steinberg module, the
network is the geometry's atlas of cube charts, the fuel is its set of magic states, and
the universal gate is the $E_6$ cubic of Section~5.}
\label{fig:stack}
\end{figure}

\subsection{The carrier: two qutrits in one photon}

A single photon has several independent degrees of freedom. Two of them --- which of three
\emph{time bins} it arrives in, and which of three \emph{frequency bins} it occupies ---
give two qutrits, and together the $\C^9$ on which the two-qutrit Paulis of Section~2 act.
Its path and polarisation give a second register, $\C^4$, whose $40$ Witting rays realise
$\W$ as \emph{states} (Theorem~5.3 and the box after it).

\begin{figure}[htbp]
\centering
\begin{tikzpicture}[>={Stealth[length=2mm]},
   el/.style={draw=#1,fill=#1!12,rounded corners=3pt,align=center,font=\sffamily\scriptsize,minimum height=11mm,minimum width=20mm}]
  \node[el=ink] (src) at (0,0) {heralded\\single photon};
  \node[el=sky] (pbs) at (2.9,0) {polarising\\beam splitter};
  \node[el=teal] (tri) at (5.8,0) {tritter $F_3$\\(time bins)};
  \node[el=gold] (cfbg) at (8.7,0) {chirped fibre\\Bragg grating};
  \node[el=rose] (det) at (11.6,0) {time- and\\frequency-\\resolved detection};
  \draw[->,thick] (src) -- (pbs); \draw[->,thick] (pbs) -- (tri);
  \draw[->,thick] (tri) -- (cfbg); \draw[->,thick] (cfbg) -- (det);
  \node[font=\sffamily\tiny,text=slate,align=center] at (2.9,-1.05) {path $\otimes$ polarisation:\\$\C^4$, Witting rays};
  \node[font=\sffamily\tiny,text=slate,align=center] at (5.8,-1.05) {qutrit Fourier gate};
  \node[font=\sffamily\tiny,text=slate,align=center] at (8.7,-1.05) {\textsc{sum}: $\ket{f,t}\mapsto\ket{f,t+f}$\\colour controls delay};
  \node[font=\sffamily\tiny,text=slate,align=center] at (11.6,-1.05) {$\C^3_{t}\otimes\C^3_{f}=\C^9$};
\end{tikzpicture}
\caption{One photon, two qutrits. The controlled-add (\textsc{sum}) gate between the
frequency and time qutrits has been demonstrated with a single passive element, a
chirped fibre Bragg grating, at basis fidelity $0.92\pm0.01$ \cite{Imany2019}.}
\label{fig:carrier}
\end{figure}

\subsection{The instruction set}

The machine's classical control state is a \emph{Pauli frame}: a label $x\in\F_3^4$ saying
which Pauli correction is pending. There are $81$ frames. Instructions are symmetries of
the geometry: the two Fourier gates, the two phase gates, the two controlled-add gates, and
the Pauli translations.

\begin{result}{Theorem 10.1 (the complete instruction group) \hfill \tierC}
The six Clifford opcodes generate exactly $\Sp(4,3)$ (order $51\,840$); adjoining one
translation generates the affine group
\[
\mathrm{ASp}(4,3)=\F_3^4\rtimes\Sp(4,3),\qquad |\mathrm{ASp}(4,3)|=81\times51\,840=4\,199\,040 .
\]
No one or two linear opcodes suffice; exactly six triples do, so \emph{four} instructions
(three Clifford, one translation) in a \emph{two-bit} opcode field generate everything, with
worst-case program length $19$ for the best triple. The minimal arithmetic core synthesises
to $43$ logic cells at $208.86$\,MHz on an iCE40 HX8K FPGA.
\cert{holonet\_machine\_blueprint\_body.tex}, Passes 2778, 2789, 2882
\end{result}

\begin{plainwords}
Four instructions are enough to perform every one of more than four million possible
frame transformations, and no program needs more than nineteen steps. The instruction set
was not designed: it is the list of symmetries of the forty-point shape, and the
theorem says which symmetries you can leave out.
\end{plainwords}

\subsection{The fuel: matter is magic}

Clifford operations alone can be simulated efficiently on an ordinary computer; quantum
advantage needs a non-Clifford resource, usually supplied as \emph{magic states}; contextuality is what
supplies them \cite{Howard2014}.

\begin{result}{Theorem 10.2 (magic is everything except one line) \hfill \tierC}
Of the $40$ Witting rays, exactly the $4$ computational-basis rays are stabiliser states;
the other $36$ are magic. The $36$ split by stabiliser fidelity into grades $8+24+4$. A
classical $\{0,1\}$ assignment can satisfy the ``exactly one'' rule on at most $36$ of the
$40$ contexts, so the contextual fraction is exactly $1/10$; the independence number of the
collinearity graph is $7$, attained by heptads of seven rays with all pairwise overlaps $1/3$.
\cert{bt822\_magic\_census\_matter\_is\_magic.py}, \cert{bt818\_ovoid\_nogo\_theta\_gap.py}
\end{result}

\begin{figure}[htbp]
\centering
\begin{tikzpicture}
  \def\r{1.55}
  \fill[ink!80] (0,0) -- (90:\r) arc[start angle=90,end angle=126,radius=\r] -- cycle;
  \fill[rose!65] (0,0) -- (126:\r) arc[start angle=126,end angle=450,radius=\r] -- cycle;
  \node[text=white,font=\scriptsize\bfseries] at (108:1.05) {$4$};
  \node[text=white,font=\small\bfseries] at (290:0.8) {$36$};
  \node[font=\sffamily\scriptsize,anchor=west,align=left] at (2.0,0.55) {\textcolor{ink}{\textbf{$4$ stabiliser rays}}: one line\\of $\W$ --- the classical readout};
  \node[font=\sffamily\scriptsize,anchor=west,align=left] at (2.0,-0.55) {\textcolor{rose!80!black}{\textbf{$36$ magic rays}}: grades $8+24+4$\\the non-classical fuel};
  \begin{scope}[xshift=8.2cm]
    \foreach \lab/\n/\c [count=\i from 0] in {{self}/1/ink,{meets $B$}/12/sky,{skew to $B$}/27/rose}{
      \node[font=\sffamily\scriptsize,anchor=east] at (0,-\i*0.55) {\lab};
      \fill[\c!65] (0.1,-\i*0.55-0.18) rectangle (0.1+\n*0.12,-\i*0.55+0.18);
      \node[font=\scriptsize,anchor=west] at (0.15+\n*0.12,-\i*0.55) {$\n$};}
    \node[font=\sffamily\tiny,text=slate,align=left,anchor=west] at (-1.3,-1.55) {contexts: $1$ classical, $12$ with one classical ray,\\$27$ fully magic --- the same $1+12+27$ as Figure~\ref{fig:bellview}};
  \end{scope}
\end{tikzpicture}
\caption{The machine's quantum resource is the geometry itself: remove one line and every
remaining point is a magic state.}
\label{fig:magic}
\end{figure}

\paragraph{The cubic gate.}
Clifford gates are ``degree two'': they act on frame labels by linear symplectic maps and
on phases by quadratic forms. The classical continuous-variable theorem of Lloyd and
Braunstein \cite{LloydBraunstein} says that quadratic operations plus a single cubic one are universal. In
Section~5 the cubic appeared as the $E_6$ form $N$, the rule by which three first-order moves
reach the top of the ladder. Pass~10947 compiles that signed cubic into a native gate on a
seven-qutrit virtual machine: ten operations in seven parallel layers, both ancillas returned
to zero. \tierC\ (\cert{w33\_pass10947\_five\_front\_execution.py}) The same object is,
in the language of Section~5, the Yukawa coupling of $E_6$ unification, and in the language of
this section, the machine's universal non-Clifford gate. \tierB

\subsection{Memory and error correction}

\begin{result}{Theorem 10.3 (the memory is the Steinberg module) \hfill \tierC}
The CSS code $[[240,81,3]]_3$ on the edges of $\W$, with distances $(d_X,d_Z)=(3,4)$,
stores its $81$ logical qutrits in $H_1(\W)$, which is the irreducible Steinberg module of
$\PSp(4,3)$. By Schur's lemma any symmetry-covariant logical error acts as a scalar.
Separately, a ternary Golay $[[11,1,5]]_3$ lane with two verified syndrome rounds has $552$
fault locations in a declared model, $2\,576\,205$ malignant fault triples, and a conservative
threshold crossing at $p\approx5.32\times10^{-4}$.
\cert{bt861\_code\_register\_is\_steinberg.py}, \cert{w33\_pass10947\_five\_front\_execution.py}
\end{result}

\subsection{The network}

\begin{figure}[htbp]
\centering
\begin{tikzpicture}[>={Stealth[length=1.8mm]},
   v/.style={circle,draw=ink!60,fill=skysoft,inner sep=0pt,minimum size=6.5mm,font=\tiny}]
  % cube chart
  \begin{scope}[scale=1.25]
    \node[v] (000) at (0,0) {000}; \node[v] (100) at (1.4,0) {100};
    \node[v] (010) at (0,1.4) {010}; \node[v] (110) at (1.4,1.4) {110};
    \node[v] (001) at (0.6,0.55) {001}; \node[v] (101) at (2.0,0.55) {101};
    \node[v] (011) at (0.6,1.95) {011}; \node[v] (111) at (2.0,1.95) {111};
    \foreach \a/\b in {000/100,000/010,100/110,010/110,001/101,001/011,101/111,011/111,000/001,100/101,010/011,110/111}{\draw[ink!55] (\a) -- (\b);}
    \draw[rose,dashed,line width=0.9pt] (000) -- (111);
    \node[font=\sffamily\scriptsize,text=slate,align=center] at (1.0,-0.65) {one chart $Q_3$: XOR routing;\\antipode $=$ collinear partner};
  \end{scope}
  % fractal
  \begin{scope}[xshift=7.2cm,yshift=1.3cm]
    \node[circle,fill=ink,text=white,minimum size=11mm,font=\scriptsize] (c) at (0,0) {$\W$};
    \foreach \k in {0,...,7}{
      \pgfmathsetmacro{\a}{\k*45}
      \node[circle,fill=teal!70,text=white,minimum size=6.5mm,font=\tiny] (s\k) at (\a:1.55) {$\W$};
      \draw[ink!40] (c) -- (s\k);
      \foreach \j in {-1,0,1}{\fill[gold] (\a+\j*14:2.25) circle (1.5pt); \draw[gold!60] (s\k) -- (\a+\j*14:2.25);}}
    \node[font=\sffamily\scriptsize,text=slate,align=center] at (0,-2.75) {recursive holonet: $40^n$ leaves,\\worst-case route $\le 8n$ reversible moves};
  \end{scope}
\end{tikzpicture}
\caption{Left: a local routing chart. Each of the $540$ pairs of disjoint lines of $\W$
carries a cube of eight points with XOR addressing. Right: substituting a copy of $\W$ at
every point gives a network with $40^n$ leaves and logarithmic routing.}
\label{fig:network}
\end{figure}

$\W$ contains exactly $540$ pairs of disjoint lines; each pair's eight points carry the
cube graph $Q_3$, in which the antipode of a vertex is its unique collinear partner. The
charts are glued along the $1620$ apartments of the Tits building of $\Sp(4,3)$. A route
needs at most $3$ in-chart moves plus $5$ chart hops, so each digit of a recursive address
costs at most $8$ reversible moves. \tierC\ (\cert{bt773\_involution\_cube\_theorem.py}, \cert{bt777\_hypercube\_chart\_web.py},
\cert{bt827\_holonet\_fractal\_architecture.py})

\begin{boundary}[What is measured, what is proved, what is not built]
\textbf{Measured in a laboratory (prior art):} the single-photon time--frequency \textsc{sum}
gate, fidelity $0.92\pm0.01$ \cite{Imany2019}. \textbf{Proved exactly:} the group theory, the
instruction semantics, reachability, code parameters, routing bounds, and FPGA synthesis
figures. \textbf{Not built:} an integrated device; a magic-state injection circuit with a
threshold; a physical demonstration of the network. The machine is a specification whose
every number can be checked, not a device.
\end{boundary}

\begin{keyidea}
The forty points are also a computer: $81$ frames, four instructions generating
$4\,199\,040$ operations, a memory that is the Steinberg module, a network of $540$ cube
charts, and a quantum fuel that is every point off one line. The $E_6$ cubic that is the
Yukawa coupling in physics is the universal gate in the machine.
\end{keyidea}
